% extensions.tex — three extension propositions of math-coding.

\section{Extensions}

Three extensions add domain-specific obligations to the foundation
schema. Each is instantiated as a packet under
\texttt{math/extensions/}.

\subsection{process-fsm: state machine with one forbidden transition}

\begin{theorem}[process-fsm.fsm]
A decision exists in one of five states
$\textsc{Draft} \mid \textsc{Applied} \mid \textsc{Reviewed} \mid \textsc{Retired} \mid \textsc{Abandoned}$.
The single forbidden transition is $\textsc{Draft} \to \textsc{Reviewed}$
without $\Witness$.
\end{theorem}

\begin{proof}
Without an explicit state machine, a packet remains ``applied'' by
default and review loses meaning. With a five-state machine, each
decision traverses an explicit path: from proposition (Draft) to
inhabitation (Applied) to human review (Reviewed), with possible
exit through Retired or Abandoned. The forbidden transition
$\textsc{Draft} \to \textsc{Reviewed}$ without $\Witness$ closes a
trivial bypass: one cannot declare a decision reviewed without
passing through inhabitation via git fix. This is the minimal
sufficient restriction: one rule closes all bypasses.
\end{proof}

\begin{theorem}[process-fsm.forbidden-transition]
$$\forall d.\ d.\mathtt{state} = \textsc{Reviewed} \;\wedge\; d.\mathtt{witness} = \bot \;\Rightarrow\; S(d) = \textsc{Fail}$$
The kernel $\mathrm{V_4}$ detects this and refuses the packet.
\end{theorem}

\subsection{dialectic-tas: thesis, antithesis, synthesis required}

\begin{theorem}[dialectic-tas.required-sections]
A packet with $\Reg = \textsc{Judgment}$ or $\Actor = \textsc{Human}$
requires the body sections \texttt{\#\# Thesis}, \texttt{\#\# Antithesis},
\texttt{\#\# Synthesis}. The structural form of dialectics is mandatory
for human judgments.
\end{theorem}

\begin{proof}
Human judgment without explicit consideration of counter-arguments
slips into self-deception. The dialectic tradition shows that
synthesis is reachable only through thesis and antithesis; collapsing
this structure turns judgment into declaration. In software
engineering, judgment about trust, security, or reversibility must
contain an explicit antithesis---the strongest counter-argument
against the decision taken. Without it, the proposition is a slogan
and the reviewer cannot distinguish considered decision from
impulsive one.
\end{proof}

\subsection{actor-discipline: signed commits fix author}

\begin{theorem}[actor-discipline.signed-commits]
A decision records its author through commit signatures. The signing
mode ($\textsc{strict} / \textsc{lenient} / \textsc{off}$) is determined
by the project through \texttt{.mathrc} and checked by the kernel.
\end{theorem}

\begin{proof}
A decision without author carries no responsibility. In distributed
development (human, agent, agentic system), commit signature is the
only cryptographically verifiable way to bind a proposition to a
subject. Without an explicit signing mode, the project ends up in a
situation where ``everyone signs, no one checks''---or vice versa,
``no one signs, authorship is untraceable''. Three modes---strict,
lenient, off---cover three classes of projects: critical
infrastructure, ordinary development, prototyping.
\end{proof}

\begin{theorem}[Signing modes]
\begin{itemize}
\item \textsc{strict}: every commit of the packet has a valid GPG/SSH
      signature. Kernel rejects the packet without signature
      ($\textsc{Fail}$).
\item \textsc{lenient}: the \texttt{amend} commit has a signature;
      others need not. Kernel returns $\textsc{Warn}$ if \texttt{amend}
      lacks signature.
\item \textsc{off}: signatures not required; kernel does not check.
\end{itemize}
\end{theorem}